\printedpage{19}
% head: A. Grothendieck : Dérivateurs, Ch. V, Édition des Universités de Montpellier II et Paris VII

in $\cM(X)$
\[
\begin{aligned}
\Ch(\cM(X)) \quad &\simeq \quad \uHom(\ \underbrace{\boldsymbol{\Phi}}_{\substack{\text{pro-category}\\ (\Phi_n)_{n\in\mathbf{N}}}}\ ,\ \underbrace{\cM(X)}_{\uHom(X,M)}\ ) \\ % reading resolved: italic M is printed; category-valued Hom underlines restored.
&\dfeq \quad \varinjlim_n \uHom(\Phi_n,\cM(X)) \\
&\simeq \quad \varinjlim_n \uHom(\Phi_n\times X,\cM) \\
&\simeq \quad \varinjlim_n \uHom(X,\uHom(\Phi_n,\cM)) \\
&\dfeq \quad \varinjlim_n \uHom(X,\Ch_n(\cM)) \\
&\simeq \quad \uHom^{!}(X,\Ch(\cM)),
\end{aligned}
\]
where the sign ! means that we take the full subcategory of functors
\[
X\longrightarrow \Ch(\cM),
\]
\msp{page 30}
which factor through one of the full subcategories $\Ch_n(\cM)$. We can also say that if we consider the direct system of categories $\Ch_n(\cM)$ as an ind-object of Cat, which we are going to denote
\[
\Chh(\cM) \qquad\qquad\qquad\qquad \text{ind-object of Cat,}
\]
we have
\[
\Ch(\cM(X)) \xrightarrow{\ \sim\ } \uHom(X,\Chh(\cM)).
\]
If we turn to $\Sigma$, $\Sigma_X$, we find at once that a path in $\cM(X)$ is a $\Sigma_X$-path, if and only if the corresponding functor $X\longrightarrow\Ch(\cM)$ factors through $\Ch_{\Sigma}(\cM)$. Thus we find
\[
\Ch_{\Sigma_X}(\cM(X))\simeq\uHom^{!}(X,\Ch_{\Sigma}(\cM)),
\]
where the sign ! has the same meaning as above. If we introduce the ind-object $\Chh_{\Sigma}(\cM)$ defined by the $\Ch_{\Sigma,n}(\cM)$, we thus also find
\[
\Ch_{\Sigma_X}(\cM(X))\xrightarrow{\ \sim\ }\uHom(X,\Chh_{\Sigma}(\cM)).
\]

\msp{page 31}
Let $c$ be a path in $\cM(X)$, and
\[
X\longrightarrow\Ch(\cM)
\]
the functor it defines, let us compose it with the source-target functors
\[
\Ch(\cM)\longrightarrow\cM\times\cM,
\]
we find a functor
\[
X\longrightarrow\cM\times\cM,
\]
i.e. an object $(F_0,F_1)$ of $\cM(X)\times\cM(X)$. I say that $F_0$ and $F_1$ are respectively the source and the target of the path $c$. A little more precisely, we have commutativity (strict) in the diagram of functors
% foot: Version 9 septembre 2001
\printedpage{20}
% head: A. Grothendieck : Dérivateurs, Ch. V, Édition des Universités de Montpellier II et Paris VII
\[
\begin{tikzcd}[column sep=large, row sep=large]
\Ch_{\Sigma_X}(\cM(X)) \arrow[r, "\sim"] \arrow[d, "{\sigma_{\cM(X)},\,\beta_{\cM(X)}}"] & \uHom^{!}(X,\Ch_{\Sigma}(\cM)) \arrow[d, "{\text{deduced from }\sigma_{\cM},\,\beta_{\cM}}"] \\
\cM(X)\times\cM(X) \arrow[r, "\sim"] & \uHom(X,\cM\times\cM),
\end{tikzcd}
\]
where the functors $\sigma$ are the source functors, the functors $\beta$ the target functors.

If we give ourselves arbitrary objects $F_0$, $F_1$ in $\cM(X)$, hence a functor
\[
X\longrightarrow\cM\times\cM,
\]
$\Ch_{\Sigma_X}(F_0,F_1)$ thus identifies with the category of ``admissible'' functors (i.e. factoring through a $\Ch_{\Sigma,n}(\cM)$) $X\xrightarrow{\ C\ }\Ch_{\Sigma}(\cM)$ which \textit{lift}
\msp{page 32 vide}
\msp{page 33}
the functor $(\sigma_{\cM},\beta_{\cM})$ \texttt{[rather} $(F_0,F_1)$\texttt{]} above, i.e. which give rise to a commutative triangle (strict)
\[
\begin{tikzcd}[row sep=large]
 & \Ch_{\Sigma}(\cM) \arrow[d, "{(\sigma_{\cM},\,\beta_{\cM})}"] \\
X \arrow[ur, "C"] \arrow[r, "{(F_0,F_1)}"] & \cM\times\cM.
\end{tikzcd}
\]
But if the functors $F_0$, $F_1$ are constant, with value $a$, $b$ as before, i.e. $(F_0,F_1)$ constant with value $(a,b)$, then the category of functors $C$ in question is tautologically isomorphic to that of functors from $X$ into the fiber of $(\sigma_{\cM},\beta_{\cM})$ at $(a,b)$, whence the fundamental isomorphism of categories
\[
\boxed{\Ch_{\Sigma_X}(X;a_X,b_X)\simeq\uHom^{!}(X,\Ch_{\Sigma}(\cM;a,b)).}
\]

We are going to make explicit the essential content of this formula in a

\textbf{Scholie.} Let $\cM$ be a category with a set $\Sigma$ of localization arrows, $a,b\in\mathrm{Ob}\,\cM$. Then we can find an object $\KK=K(a,b)$, $K=(K_n)$, of Ind\,Cat (indexed by $\mathbf{N}$ but this fact matters little), and a canonical isomorphism for every category $X$ (isomorphism functorial in $X$)
\[
\begin{aligned}
\mathrm{Hom}_{\Sigma_X}(a_X,b_X) \quad &\simeq \quad \pi_0(\uHom(X,\KK)) \\
&\dfeq \quad \pi_0(\ \varinjlim \uHom(X,K_n)) \\
&\simeq \quad \varinjlim \pi_0(\uHom(X,K_n)).
\end{aligned}
\]
% foot: Version 9 septembre 2001

\printedpage{21}
% head: A. Grothendieck : Dérivateurs, Ch. V, Édition des Universités de Montpellier II et Paris VII

\msp{page 34}
I am going to dwell on this result, without worrying for the moment about the particular way in which the ind-category $K$ \texttt{[rather} $\KK$\texttt{]} was obtained in terms of $\cM$, $\Sigma$, $a$, $b$.

Since the $\pi_0\uHom(X,K_n)=\mathrm{Hom}_{\Cathot}(X,K_n)$ depend only on $X$ and on $K_n$ as objects of Cat\,hot, I am going to interpret $\KK$ as an ind-object of Cat\,hot, and write the fundamental formula in the form
\[
\mathrm{Hom}_{\Sigma_X}(a_X,b_X)\simeq\mathrm{Hom}_{\mathrm{Ind\,Cat\,hot}}(X,\KK).
\]

Thus, the statement of the scholie implies two things:

\begin{enumerate}
\item[a)] The contrafunctor $X\longmapsto\mathrm{Hom}_{\Sigma_X}(a_X,b_X)$ on Cat factors canonically through Cat\,hot, via a canonical functor
\[
\Cathot^{o}\ \xrightarrow{\ H=H^{\cM,\Sigma}_{a,b}\ }\mathrm{Ens}\ ;
\]
\item[b)] This latter contrafunctor on Cat\,hot is \textit{ind-representable} by an ind-object $\KK$ of Cat\,hot.
\end{enumerate}

It is this ind-object of Cat\,hot which is defined up to unique isomorphism by $\cM$, $\Sigma$, $a$, $b$, independently of any computational artifice for making the $\mathrm{Hom}_{\Sigma_X}(a_X,b_X)$ explicit.

Let us also note the following conclusion, by no means obvious a priori:

\msp{page 35}
Let us first note (this being implicit in the scholie) that every arrow in Cat
\[
X\xrightarrow{\ f\ }Y
\]
defines
\[
\cM(Y)\xrightarrow{\cM(f)}\cM(X) \qquad\qquad\qquad \text{mapping } \Sigma_Y \text{ into } \Sigma_X,
\]
whence
\[
\cM(Y)\Sigma_Y^{-1}\dfeq\mathbf{D}_{\cM,\Sigma}(Y)\longrightarrow\cM_X\Sigma_X^{-1}=\mathbf{D}_{\cM,\Sigma}(X),
\]
\texttt{[rather} $\cM(X)$\texttt{]} whence, for every $a,b\in\mathrm{Ob}\,\cM$, since $\cM(f)$ transforms $a_Y$, $b_Y$ into $a_X$, $b_X$:
\[
\boxed{\mathrm{Hom}_{\Sigma_Y}(a_Y,b_Y)\longrightarrow\mathrm{Hom}_{\Sigma_X}(a_X,b_X).}
\]

This being so, we have:

\textbf{Corollary:} \textit{If $f$ is a homotopism, then for any $a,b\in\cM$, the preceding arrow is an isomorphism.} % [sic] quelque soient -> quels que soient, grammatical form of the locution

The natural question is to know whether the same conclusion is still valid, assuming only that $f$ is a ``\textit{quasi-isomorphism}'' (or ``weak equivalence''), here that $f$ defines an isomorphism in the category of fractions Hot of Cat. (Not to be confused with Cat\,hot,
% foot: Version 9 septembre 2001
\printedpage{22}
% head: A. Grothendieck : Dérivateurs, Ch. V, Édition des Universités de Montpellier II et Paris VII
\msp{page 36}
distinctly finer). If this were so, the contrafunctor
\[
X\longmapsto\mathrm{Hom}_{\Sigma_X}(a_X,b_X)
\]
would factor not only through Cat\,hot, but even through the coarser category of fractions Hot. And then the question would arise whether this hypothetical contrafunctor
\[
\mathrm{Hot}^{o}\longrightarrow\mathrm{Ens}
\]
is still ind-representable; and preferably (\footnote[5]{this will be automatic if we have ind-representability as a functor on Hot, cf. below.}) ind-representable by the ind-object image of $\KK$ in Ind\,Hot under
\[
\mathrm{Ind}(\Cathot)\longrightarrow\mathrm{Ind}(\mathrm{Hot}),
\]
deduced from the localization functor $\Cathot\longrightarrow\mathrm{Hot}$. (For this ind-object is a ready-made candidate.)

\msp{page 37}
I am going to give a few categorical reminders, generalizing the situation of the canonical localization functor
\[
\Cathot\xrightarrow{\ \varphi\ }\mathrm{Hot},
\]
which we know admits a right adjoint \texttt{[if a reader has a proof of this fact, please contact G. Maltsiniotis: [e-mail address removed]]}
\[
\Cathot\xleftarrow{\ \psi\ }\mathrm{Hot}, \qquad\qquad\qquad \mathrm{Hom}_{\mathrm{Hot}}(\varphi X,Y)\simeq\mathrm{Hom}_{\Cathot}(X,\psi Y),
\]
where $\psi$ is the ``Kan envelope'' or ``Kan representative'' functor. We are therefore going to consider the more general situation
\[
\begin{tikzcd}[column sep=large]
C \arrow[r, shift right=1, "\varphi"'] & C' \arrow[l, shift right=1, "\psi"']
\end{tikzcd}
\qquad\qquad \mathrm{Hom}_{C'}(\varphi X,X')\simeq\mathrm{Hom}_{C}(X,\psi X')
\]
i.e.
\[
\varphi,\psi \text{ pair of adjoint functors.}
\]
We assume moreover
\[
\left\{
\begin{aligned}
&\varphi \text{ a localization functor, i.e. } \psi \text{ fully faithful,}\\
&\text{i.e. } \varphi\psi\xrightarrow[\sim]{\text{adj}}\mathrm{id}_{C'} \quad .
\end{aligned}
\right.
\]
Let $W=W_{\varphi}\subset\mathrm{Fl}\,C$, defined by $\varphi$: $W_{\varphi}=\{u\in\mathrm{Fl}\,C\mid\varphi(u)\text{ isomorphism}\}$. Here is the proposition I had forgotten, and which I had trouble reconstructing, alas!

\textbf{Proposition 1.} \textit{Let $K\in \mathit{Ob}C$. The following conditions on $K$ are equivalent:}
% foot: Version 9 septembre 2001

\printedpage{23}
% head: A. Grothendieck : Dérivateurs, Ch. V, Édition des Universités de Montpellier II et Paris VII

\begin{enumerate}
\item[(i)] \textit{The functor}
\[
X\longmapsto K(X)\dfeq\mathrm{Hom}_C(X,K) \qquad\quad C^{o}\longrightarrow\mathrm{Ens}
\]
\textit{transforms arrows of $W$ (``quasi-isomorphisms'') into bijections.}
\item[(i$'$)] \textit{Same condition for the adjunction arrows}
\[
X\longrightarrow\psi\varphi X \qquad\quad (\forall X\in\mathrm{Ob}\ \mathrm{C})
\]
\item[(i$''$)] \textit{Same condition for the adjunction arrow} (\footnote[6]{N.B. It even suffices that the element $\mathrm{id}_K$ of $F_K(K)$ come from $F_K(\psi\varphi K)$ \texttt{[where} $F_K=K(-)=\mathrm{Hom}(-,K)$\texttt{]}, i.e.\par (i$'''$) There exists a retraction of $\psi\varphi X$ onto $X$ \texttt{[rather} $\psi\varphi K$ \texttt{onto} $K$\texttt{]} (NB This will even be an isomorphism).})
\[
K\longrightarrow\psi\varphi K.
\]
\item[(ii)] \textit{$K$ belongs to the essential image of $\psi$}
\item[(ii$'$)] \textit{$K\longrightarrow\psi\varphi K$ (adjunction) is an isomorphism.}
\item[(ii$''$)] \textit{The homomorphism functorial in $X$}
\[
\mathrm{Hom}_C(X,K)\longrightarrow\mathrm{Hom}_{C'}(\varphi X,\varphi K)
\]
\textit{is an isomorphism.}
\end{enumerate}

\textsc{Proof.} We first have as quasi-tautologies (ii$'$)$\Longrightarrow$ (ii)$\Longrightarrow$ (ii$''$)$\Longrightarrow$ (ii$'$), i.e. the equivalence of the last three conditions. Indeed, (ii$'$)$\Longrightarrow$ (ii) tautological, (ii)$\Longrightarrow$ (ii$''$) since if $K=\psi(K')$ we have % reading resolved: single and double mathematical primes as printed.
\msp{page 38}
by adjunction
\[
\mathrm{Hom}_C(X,K)=\mathrm{Hom}_C(X,\psi(K'))\overset{\text{adj}}{\simeq}\mathrm{Hom}_{C'}(\varphi X,K'),
\]
now $\varphi(K)=\varphi\psi(K')\xrightarrow{\ \sim\ }K'$ (adjunction, and $\psi$ fully faithful).
\[
\begin{tikzcd}[column sep=small, row sep=large]
\mathrm{Hom}_C(X,K) \arrow[rr, "\sim"] \arrow[dr, "\varphi_{X,K}"] & & \mathrm{Hom}_C(\varphi X,K') & \text{\footnotesize which is completed to} \\ % [sic] Hom_C(\varphi X,K') -> Hom_{C'}(\varphi X,K'), \varphi X and K' are objects of C' (cf. the display above)
& \mathrm{Hom}_{C'}(\varphi X,\ \underbrace{\varphi K}_{=\varphi\psi(K')}\ ) \arrow[ur, "\sim", "\text{adj}"'] & &
\end{tikzcd}
\]
which proves that $\varphi_{X,K}$ is an isomorphism.

(ii$''$) $\Longrightarrow$ (ii$'$). For the arrow
% foot: Version 9 septembre 2001